% ============================================================
%  NF Validation Research Strategies — Main LaTeX Document
%  Author: Alok Prakash
%  This document synthesizes research findings for improving
%  Yaksha-Prashna in NF Validation and Network Verification
% ============================================================

\documentclass[12pt, a4paper]{article}

% ---------- Packages ----------
\usepackage[margin=1in]{geometry}
\usepackage[T1]{fontenc}
\usepackage[utf8]{inputenc}
\usepackage{lmodern}
\usepackage{microtype}
\usepackage{hyperref}
\usepackage{xcolor}
\usepackage{graphicx}
\usepackage{booktabs}
\usepackage{longtable}
\usepackage{array}
\usepackage{enumitem}
\usepackage{amsmath}
\usepackage{amssymb}
\usepackage{mdframed}
\usepackage{tcolorbox}
\usepackage{titlesec}
\usepackage{fancyhdr}
\usepackage{setspace}
\usepackage{soul}
\usepackage{listings}
\usepackage{tocloft}
\usepackage[normalem]{ulem}

% ---------- Colors ----------
\definecolor{ypblue}{RGB}{30, 90, 160}
\definecolor{strategygreen}{RGB}{34, 120, 56}
\definecolor{warningorange}{RGB}{200, 100, 20}
\definecolor{lightgray}{RGB}{245, 245, 245}
\definecolor{darkgray}{RGB}{60, 60, 60}
\definecolor{accentred}{RGB}{180, 30, 30}

% ---------- Custom Environments ----------
\tcbuselibrary{skins, breakable}

\newtcolorbox{strategybox}{
  enhanced,
  breakable,
  colback=strategygreen!8,
  colframe=strategygreen!70!black,
  title=\textbf{Research Strategy for Yaksha-Prashna},
  fonttitle=\small\bfseries,
  attach boxed title to top left={yshift=-2mm, xshift=4mm},
  boxed title style={colback=strategygreen!80!black, colframe=strategygreen!80!black},
  left=4pt, right=4pt, top=4pt, bottom=4pt
}

\newtcolorbox{analogybox}{
  enhanced,
  breakable,
  colback=ypblue!7,
  colframe=ypblue!60,
  title=\textbf{Analogy to Yaksha-Prashna},
  fonttitle=\small\bfseries,
  attach boxed title to top left={yshift=-2mm, xshift=4mm},
  boxed title style={colback=ypblue!80, colframe=ypblue!80},
  left=4pt, right=4pt, top=4pt, bottom=4pt
}

\newtcolorbox{limitbox}{
  enhanced,
  breakable,
  colback=warningorange!8,
  colframe=warningorange!70,
  title=\textbf{Limitations w.r.t.\ Yaksha-Prashna},
  fonttitle=\small\bfseries,
  attach boxed title to top left={yshift=-2mm, xshift=4mm},
  boxed title style={colback=warningorange!80, colframe=warningorange!80},
  left=4pt, right=4pt, top=4pt, bottom=4pt
}

\newtcolorbox{insightbox}{
  enhanced,
  breakable,
  colback=lightgray,
  colframe=darkgray!50,
  title=\textbf{Key Insight},
  fonttitle=\small\bfseries,
  left=4pt, right=4pt, top=4pt, bottom=4pt
}

% ---------- Header/Footer ----------
\pagestyle{fancy}
\fancyhf{}
\fancyhead[L]{\small\textcolor{darkgray}{NF Validation Research Strategies}}
\fancyhead[R]{\small\textcolor{darkgray}{Alok Prakash}}
\fancyfoot[C]{\small\thepage}
\renewcommand{\headrulewidth}{0.4pt}

% ---------- Section Formatting ----------
\titleformat{\section}[block]
  {\Large\bfseries\color{ypblue}}
  {\thesection}{1em}{}[\titlerule]

\titleformat{\subsection}[block]
  {\large\bfseries\color{darkgray}}
  {\thesubsection}{1em}{}

\titleformat{\subsubsection}[runin]
  {\normalsize\bfseries\color{strategygreen}}
  {\thesubsubsection}{1em}{}[.]

% ---------- Hyperref Setup ----------
\hypersetup{
  colorlinks=true,
  linkcolor=ypblue,
  citecolor=strategygreen,
  urlcolor=accentred,
  pdftitle={NF Validation Research Strategies},
  pdfauthor={Alok Prakash},
}

% ---------- Listings ----------
\lstset{
  basicstyle=\small\ttfamily,
  backgroundcolor=\color{lightgray},
  frame=single,
  framesep=4pt,
  rulecolor=\color{darkgray!40},
  breaklines=true,
}

% ============================================================
\begin{document}
% ============================================================

% ---------- Title Page ----------
\begin{titlepage}
  \centering
  \vspace*{2cm}
  {\Huge\bfseries\color{ypblue} NF Validation Research Strategies\par}
  \vspace{0.5cm}
  {\large\color{darkgray} Findings from Selected Papers for Improving Yaksha-Prashna\\
  in NF Validation and Network Verification\par}
  \vspace{2cm}
  {\large\bfseries Alok Prakash\par}
  \vspace{0.5cm}
  {\normalsize IIT Hyderabad \par}
  \vspace{0.3cm}
  {\normalsize \today\par}
  \vspace{3cm}
  \begin{mdframed}[backgroundcolor=lightgray, linecolor=ypblue, linewidth=1.5pt]
    \centering
    \small\textit{This document is a living research strategy record synthesizing insights
    from a curated literature survey of 89 papers (filtered to 16 high-impact works)
    on NF validation and network verification, with the explicit goal of identifying
    concrete improvements for Yaksha-Prashna.}
  \end{mdframed}
\end{titlepage}

% ---------- Table of Contents ----------
\tableofcontents
\newpage

% ============================================================
\section{Problem Statement and Research Context}
\label{sec:problem}
% ============================================================

\subsection{What is Yaksha-Prashna?}

Yaksha-Prashna is a pioneering system for functional correctness verification of eBPF Network Functions (NFs) directly from compiled bytecode. Among its notable achievements, it operates on raw eBPF XDP/TC bytecode without requiring source code, and it performs multi-program chain verification — two capabilities unmatched by prior work in the field.

Technically, Yaksha-Prashna operates as follows:
\begin{enumerate}
  \item \textbf{Input}: Compiled eBPF object file (ELF binary) — no source code, no annotations required.
  \item \textbf{CFG-NC Construction}: Builds a \emph{Control Flow Graph with Network Context}, capturing:
    \begin{itemize}
      \item Packet header field accesses (which fields are read/written)
      \item BPF map read/write operations (which maps, what keys)
      \item Helper function call sequences (\texttt{bpf\_map\_lookup\_elem}, \texttt{bpf\_redirect}, etc.)
      \item XDP/TC action decisions (\texttt{XDP\_PASS}, \texttt{XDP\_DROP}, \texttt{TC\_ACT\_OK}, etc.)
    \end{itemize}
  \item \textbf{Prolog Query Engine}: A Prolog-based query engine evaluates behavioral assertions against the CFG-NC, answering questions like ``Does this program ever drop packets matching predicate $P$?'' or ``Does NF$_i$'s output ever reach NF$_{i+1}$ with field $F$ unmodified?''
  \item \textbf{Chain Analysis}: Handles eBPF NF chains by composing CFG-NC representations across multiple programs, tracking stateful carry-over via shared BPF map dependency analysis.
\end{enumerate}

\subsection{Current Limitations}

Despite these contributions, Yaksha-Prashna has known limitations that motivate this research survey:

\begin{enumerate}
  \item \textbf{Limited stateful field tracking}: Only 6 packet header fields are modeled in the stateful component of CFG-NC analysis, limiting the scope of behavioral queries.
  \item \textbf{Path explosion risk}: Complex NFs with large state spaces can exhaust the static analysis budget, a problem shared by symbolic execution approaches.
  \item \textbf{Narrow behavioral query vocabulary}: The Prolog query language lacks temporal operators (no LTL-style ``eventually'' or ``always before''), limiting expressiveness for multi-packet chain properties.
  \item \textbf{No formal BPF map composition model}: While YP tracks BPF map dependencies across chains, there is no formal algebraic model of map state composition.
  \item \textbf{No performance-level assertions}: YP verifies functional correctness but cannot assert performance properties (e.g., ``this NF processes $\geq X$ Mpps for $\leq 64$B packets'').
  \item \textbf{No probabilistic or approximate reasoning}: All assertions are binary (holds/violates) with no gradation.
\end{enumerate}

\subsection{Research Agenda}

This document records our research strategy as we survey the literature to find concrete, actionable improvements for Yaksha-Prashna. We have surveyed 89 papers in the domain of NF validation, network verification, and eBPF program analysis. After filtering for relevance, methodological novelty, and compatibility with Yaksha-Prashna's bytecode-level approach, we have identified \textbf{16 high-impact papers}. This document contains in-depth analyses of 9 of these, each analyzed through the lens of: ``What can YP concretely borrow from this work?''

The spreadsheet tracking all papers is maintained at:\\
\url{https://docs.google.com/spreadsheets/d/1co271Fl6_rgMSXye8-byWxMWKNp_goEyqrY8yG3DUQo/edit}

% ============================================================
\section{Research Analysis: Heimdall}
\label{sec:heimdall}
% ============================================================

\subsection*{Paper: Heimdall: LLM + Symbolic Execution for eBPF C-to-Rust Migration Verification}
\textbf{Year}: 2026 (preprint) \quad \textbf{arXiv}: 2605.25411 \quad \textbf{Authors}: Vasu Dasu, Monika Sharma et al., Penn State / UIC

\subsubsection{First-Pass Intuition}

This paper addresses a specific and dangerous class of bugs in eBPF programs: \emph{source-level bugs that corrupt data and leak previously traced events to userspace}. Using this leaked information, an attacker can recover data normally hidden or secured by KASLR (Kernel Address Space Layout Randomization --- a cybersecurity mitigation that randomizes kernel memory locations). The authors discovered unreported leaks in 10 open-source eBPF programs whose ring buffer or stack records contained previously traced data.

Heimdall is an automated pipeline to harden such verifier-accepted but buggy programs. The eBPF verifier accepts these programs because they are \emph{safe} by the verifier's definition (no out-of-bounds memory access, no null pointer dereferences) --- but safety is not the same as security. A program can be safe yet leak sensitive kernel memory to userspace through poorly managed ring buffer residue.

\subsubsection{Classification as NF Validation}

\begin{insightbox}
This paper is \emph{borderline} NF validation. Its primary contribution is \textbf{compiler verification / security hardening}: it validates that a migrated eBPF program does not introduce new information leaks. However, at its core, it \emph{is} validating an eBPF program's behavioral equivalence --- which is a form of NF behavioral validation. Most importantly, the methodology (symbolic execution on eBPF bytecode + Z3 equivalence) is directly applicable to NF behavioral validation even if that was not the authors' primary intent.
\end{insightbox}

\subsubsection{Technical Mechanism}

Heimdall is a 5-stage automated pipeline for C$\to$Rust eBPF migration with formal equivalence verification:

\begin{enumerate}
  \item \textbf{LLM Translation}: A GPT-class LLM translates legacy eBPF C code to Rust-eBPF (Aya framework).
  \item \textbf{Compile + KV Verify}: Both C and Rust programs are compiled to eBPF bytecode. Key-value pair analysis checks basic structural equivalence.
  \item \textbf{Static Safety Engine}: A rule-based static checker verifies that the Rust translation does not introduce new unsafe patterns.
  \item \textbf{Symbolic Execution (angr backend)}: A custom eBPF backend for the \texttt{angr} Python binary analysis framework symbolically executes both the original C-compiled bytecode and the Rust-compiled bytecode. This is a novel contribution: the team built a full eBPF VEX IR lifter, ELF loader, 20+ BPF map type models, and all BPF helpers as symbolic stubs.
  \item \textbf{Z3 Equivalence Check + CEGIS Repair}: Z3 is invoked to prove that for all inputs, both programs produce identical packet decisions and map state changes. When Z3 finds a counterexample, a CEGIS (Counterexample-Guided Inductive Synthesis) loop re-prompts the LLM for a corrected translation.
\end{enumerate}

\textbf{Key tools}: \texttt{angr} (Python binary analysis), Z3 SMT solver, GPT-class LLM (for translation), Clang/LLVM (for eBPF compilation), Aya (Rust eBPF framework).

\textbf{Novel data structure}: \emph{ITE-chain map encoding} --- bitvector-based encoding of BPF map semantics as if-then-else chains in Z3, replacing Theory of Arrays (which is incompatible with angr/Claripy).

\textbf{Achieved}: 96/102 (94.1\%) formally proven-equivalent translations on real-world eBPF programs.

\subsubsection{Security Findings}

Beyond NF migration, Heimdall discovered real security vulnerabilities:
\begin{itemize}
  \item \textbf{KASLR-defeating stack leak} in \texttt{bashreadline}
  \item \textbf{9 cross-event ring-buffer residue leaks} in \texttt{mountsnoop}, \texttt{opensnoop}, and other bpftrace tools
\end{itemize}
These were previously verifier-accepted programs --- demonstrating that the eBPF verifier's safety guarantees are insufficient for security.

\begin{analogybox}
\textbf{What Heimdall does with angr + Z3} is directly analogous to \textbf{what YP does with CFG-NC + Prolog}:
\begin{itemize}
  \item Both operate on eBPF bytecode (no source code required)
  \item Both extract a behavioral model from the bytecode (Heimdall: symbolic states; YP: CFG-NC dataflow)
  \item Both query properties over that model (Heimdall: Z3 equivalence; YP: Prolog assertions)
  \item Heimdall's ITE-chain BPF map encoding is the symbolic-execution counterpart of YP's map dependency tracking
\end{itemize}
The critical difference: Heimdall does \emph{equivalence checking} (are two programs behaviorally identical?), while YP does \emph{assertion checking} (does this program satisfy property $P$?). Both are needed for full NF validation.
\end{analogybox}

\begin{strategybox}
\textbf{Strategy 1 — Adopt angr's eBPF backend as a symbolic layer for YP}:\\
Heimdall's custom angr eBPF backend (20+ map type models, all helper stubs, VEX IR lifter) represents significant engineering work for symbolically modeling eBPF behavior. YP could integrate this as an optional \emph{symbolic enrichment layer}: when Prolog queries require deeper path-level reasoning than static dataflow can provide, fall through to angr-based symbolic execution for targeted paths.

\textbf{Strategy 2 — Adapt ITE-chain map encoding for YP's BPF map modeling}:\\
YP currently tracks BPF map operations (read/write) at the dataflow level. Heimdall's ITE-chain encoding provides a richer mathematical model of map semantics that preserves last-write-wins ordering under symbolic key aliasing. This can extend YP's map analysis to handle aliased key accesses --- a current limitation.

\textbf{Strategy 3 — Use CEGIS loop as query refinement}:\\
When a YP Prolog query is too coarse-grained and produces spurious results, a CEGIS-style loop (counterexample from Z3 $\to$ refine Prolog clause $\to$ re-verify) could improve assertion precision.

\textbf{Strategy 4 — Semantic enrichment, not replacement}:\\
Heimdall should be used as a \emph{semantic enrichment layer} for YP: run Heimdall to detect information leaks and safety issues, then YP for functional behavioral assertions. The two are complementary, not competing.
\end{strategybox}

\begin{limitbox}
\begin{itemize}
  \item Heimdall requires C \emph{source code} as input (for LLM translation). YP needs no source --- this is YP's advantage.
  \item Heimdall is single-program only: no NF chain analysis, no inter-program reasoning.
  \item Symbolic execution (angr) faces path explosion for complex NFs --- Heimdall acknowledges this limitation.
  \item Heimdall's equivalence checking proves $A = B$ for two programs, not $A \models P$ for a property $P$ --- it cannot be used directly for behavioral assertion queries.
\end{itemize}
\end{limitbox}

% ============================================================
%  Sections inserted by research analyst subagents
% ============================================================

% ============================================================
%  Section: VigNAT and Vigor Analysis
% ============================================================

\section{Research Analysis: VigNAT and Vigor}
\label{sec:vignat-vigor}

\subsection*{Papers}
\begin{itemize}
  \item \textbf{VigNAT}: \emph{A Formally Verified NAT}, A.\ Zaostrovnykh, S.\ Pirelli, L.\ Pedrosa, K.\ Argyraki, G.\ Candea. ACM SIGCOMM 2017.
  \item \textbf{Vigor}: \emph{Verifying Software Network Functions with No Verification Expertise}, A.\ Zaostrovnykh et al. ACM SOSP 2019.
\end{itemize}

These two papers from EPFL's Dependable Systems Lab represent the most complete vision of \emph{end-to-end formal NF verification} prior to the eBPF era. VigNAT proves a single NAT implementation correct; Vigor generalizes the framework to arbitrary NFs. Together, they constitute the gold standard against which all subsequent NF verification work is measured --- including Yaksha-Prashna.

% ------------------------------------------------------------
\subsection{VigNAT: A Formally Verified NAT}
% ------------------------------------------------------------

\subsubsection{First-Pass Intuition}

VigNAT is a Network Address Translator (NAT) written in C that has been \emph{mathematically proven} to be: (1) crash-free, (2) memory-safe, and (3) semantically correct according to RFC 3022 (the NAT specification standard) --- for \emph{every possible sequence of packets, forever}.

This is a remarkable achievement. Most network engineers would test a NAT by throwing traffic at it and checking outputs. VigNAT instead constructs a \emph{mathematical proof} that the NAT is correct for \emph{all possible inputs}, including adversarial packet sequences, edge cases in port allocation, and ICMP error handling.

The central challenge they address is \textbf{path explosion in symbolic execution}. A NAT with connection state has an enormous state space: every possible combination of active connections, port mappings, and timeout states creates a distinct execution path. Symbolic execution tools like KLEE would explode trying to explore all of them simultaneously.

\begin{insightbox}
\textbf{The key insight of VigNAT}: Decompose the NF into two parts with different verification strategies, then combine the results using \emph{lazy proofs}.
\begin{enumerate}
  \item \textbf{Data structure library} (\texttt{libVig}): the stateful component --- hash maps, vectors, port allocators. This is the hard part. Prove it correct using \emph{VeriFast separation logic} (interactive formal proofs with heap reasoning).
  \item \textbf{Stateless NAT logic}: the packet-processing code that uses the data structures. This is the easy part. Verify it using \emph{KLEE symbolic execution} --- no path explosion because state is abstracted away by the verified library.
\end{enumerate}
\end{insightbox}

\subsubsection{Technical Mechanism}

\paragraph{Step 1 --- Formalize the RFC.}
RFC 3022 defines correct NAT behavior in English prose. The VigNAT team translated this into a formal Python specification: a mathematical state machine with precise preconditions and postconditions for every packet type (TCP, UDP, ICMP).

\paragraph{Step 2 --- Implement libVig with annotations.}
Write the NAT data structures (Map, Vector, DoubleMap, PortAllocator) in C with VeriFast separation logic annotations. These annotations describe memory ownership (\emph{who owns which heap block}), invariants (\emph{the map always has a valid length}), and preconditions/postconditions for each function.

\paragraph{Step 3 --- Prove libVig correct (VeriFast).}
Run VeriFast on the annotated C code. VeriFast statically checks that the separation logic annotations are consistent and that the implementation correctly implements the specified behavior. This produces a \emph{machine-checked proof} that \texttt{libVig} is correct.

\paragraph{Step 4 --- Verify NAT logic (KLEE).}
The NAT packet-processing logic is now written using only \texttt{libVig} API calls --- no raw pointer manipulation. KLEE symbolically executes this logic with symbolic inputs (all possible packet headers, all possible prior states). Because \texttt{libVig} operations are modeled as abstract operations (not concrete implementations), KLEE's state space is manageable.

\paragraph{Step 5 --- Assert RFC compliance.}
KLEE checks that for every symbolic execution path, the NAT's output satisfies the RFC 3022 formal specification. If KLEE finds a violation, it generates a concrete packet sequence that triggers the bug.

\paragraph{Step 6 --- Combine with lazy proofs.}
VeriFast proves \texttt{libVig} correct; KLEE proves the NAT logic correct \emph{assuming} a correct library. By composing these proofs, the full system is proven correct. The ``lazy'' aspect: library correctness is proven once and reused across all NFs that use \texttt{libVig}.

\textbf{Tools}: VeriFast, KLEE, Python (RFC spec), DPDK (runtime), Z3 (implicit, via KLEE), Clang/LLVM.

\textbf{Performance}: 1.8 Mpps throughput, $\sim$0.38 $\mu$s per-packet latency --- competitive with unverified DPDK NATs.

\subsubsection{What VigNAT Proves}

\begin{enumerate}
  \item RFC 3022 port mapping correctness (endpoint-independent and endpoint-dependent)
  \item Session consistency (every connection has a consistent bi-directional mapping)
  \item Port allocation correctness (no double-allocation, no port exhaustion silent failure)
  \item ICMP error translation correctness
  \item Session timeout behavior
  \item Memory safety (no buffer overflows, no use-after-free)
  \item Crash freedom (no assertions fail, no undefined behavior)
\end{enumerate}

\begin{analogybox}
\textbf{VigNAT's DS split $\leftrightarrow$ YP's BPF map abstraction}

VigNAT's most important architectural decision --- splitting the NF into ``hard stateful DS'' (formally proven) + ``stateless logic'' (symbolically executed) --- is precisely the structure YP needs for its BPF map analysis:

\begin{itemize}
  \item VigNAT's \texttt{libVig} $\leftrightarrow$ YP's BPF map dependency tracking
  \item VigNAT's KLEE over stateless logic $\leftrightarrow$ YP's Prolog queries over CFG-NC
  \item VigNAT's lazy proofs $\leftrightarrow$ YP's composition of per-NF analyses in chains
\end{itemize}

The key difference: VigNAT requires \emph{3,000+ lines of manual VeriFast annotations}. YP extracts equivalent structural information automatically from the bytecode, with zero developer burden. This is YP's core advantage.
\end{analogybox}

\begin{strategybox}
\textbf{Strategy A --- Adopt the ``DS split'' for YP's BPF map analysis}:\\
Rather than treating all BPF map operations uniformly, classify them:
\begin{enumerate}
  \item \emph{Complex stateful maps} (hash maps with aliased keys, LRU maps): model formally using ghost map axioms (from Klint/VigNAT inspiration).
  \item \emph{Simple array maps} (per-CPU counters, configuration arrays): analyze with lightweight dataflow, no formal model needed.
\end{enumerate}
This reduces analysis complexity for simple NFs while maintaining soundness for complex ones.

\textbf{Strategy B --- Use VigNAT's RFC formalization as a query template}:\\
VigNAT's Python formalization of RFC 3022 can be adapted as a Prolog query template for YP. When a user asks ``does this NAT implement RFC 3022?'', YP would run the RFC 3022 Prolog assertions against the bytecode's CFG-NC.

\textbf{Strategy C --- Lazy proof composition for chain verification}:\\
VigNAT's lazy proof idea (prove library once, reuse across NFs) directly applies to YP's NF chain analysis: compute per-NF behavioral summaries once, then compose them for chain-level queries without re-analyzing individual NFs.
\end{strategybox}

\begin{limitbox}
\begin{itemize}
  \item VigNAT requires \textbf{C source code} with \textbf{3,000+ lines of manual VeriFast annotations}. YP requires neither.
  \item VigNAT is \textbf{single-NF only}: no chain verification, no multi-NF composition.
  \item VigNAT is specific to \textbf{DPDK-based NFs}. YP targets eBPF/XDP, a different (and newer) ecosystem.
  \item VigNAT's verification is \textbf{not push-button}: it requires significant verification expertise. Vigor addresses this.
\end{itemize}
\end{limitbox}

% ------------------------------------------------------------
\subsection{Vigor: Verifying Software NFs with No Verification Expertise}
% ------------------------------------------------------------

\subsubsection{First-Pass Intuition}

Vigor is the natural successor to VigNAT: it generalizes VigNAT's approach into a \emph{framework} that any developer can use without verification expertise. The vision is compelling: an NF developer writes their NF in C using Vigor's library, runs a single command, and gets a formal proof of correctness back.

This is the ``push-button'' promise. The developer does not write VeriFast annotations, does not understand separation logic, and does not configure KLEE. The verification is fully automated --- the framework handles all of it.

Vigor extends VigNAT in two key ways:
\begin{enumerate}
  \item \textbf{Broader NF support}: Not just NAT, but also stateful firewall, load balancer, and others.
  \item \textbf{Composable verification}: Each NF component is verified independently and composed, making the framework extensible.
\end{enumerate}

\subsubsection{Technical Mechanism}

\paragraph{Framework Architecture.}
Vigor provides a \textbf{Vigor library} (Vigor STL) of verified data structures built on \texttt{libVig}. Developers write NFs using only this library. The framework then automatically:

\begin{enumerate}
  \item \textbf{Lifts the NF to a formal model}: Vigor's front-end uses Clang to parse the NF C code and extract a formal model of its behavior.
  \item \textbf{Runs KLEE with Vigor stubs}: The formal model is symbolically executed with Vigor library functions replaced by their formal specifications (stubs that capture the abstract behavior).
  \item \textbf{Generates a Python spec}: Vigor's toolchain produces a Python behavioral specification of the NF automatically from the implementation.
  \item \textbf{Checks spec compliance}: KLEE verifies that the NF implementation satisfies its auto-generated spec for all inputs.
  \item \textbf{Push-button result}: Developer gets: verified / not verified (with counterexample).
\end{enumerate}

\paragraph{Composability.}
Vigor's key architectural insight is \emph{modular verification}: each Vigor library component is verified once, and its verification certificate is reused across all NFs. This is the generalization of VigNAT's ``lazy proofs'' to an entire ecosystem of verified components.

\textbf{Tools}: Clang/LLVM, KLEE, VeriFast (for library components), Python (spec generation), DPDK.

\textbf{NFs verified}: NAT, Firewall, Load Balancer (all DPDK-based).

\begin{analogybox}
\textbf{Vigor's push-button framework $\leftrightarrow$ YP's zero-expertise bytecode analysis}

Vigor and YP share the same goal: \emph{make formal NF verification accessible without requiring users to be verification experts}. The approaches differ fundamentally:

\begin{itemize}
  \item Vigor achieves this by \emph{constraining the developer}: use our library, follow our structure.
  \item YP achieves this by \emph{operating after the fact}: analyze the compiled bytecode of any eBPF NF, regardless of how it was written.
\end{itemize}

YP's approach is strictly more general: it does not constrain the developer's toolchain, language, or library choices. This is a critical positioning argument.

Vigor's composable verification certificates are the conceptual ancestor of YP's per-NF CFG-NC summaries that are composed for chain analysis. The difference: Vigor's certificates are manual (library annotations); YP's summaries are automatically extracted from bytecode.
\end{analogybox}

\begin{strategybox}
\textbf{Strategy D --- Design YP as a framework, not just a tool}:\\
Vigor's success comes from being a \emph{framework} that developers build within. YP should similarly provide:
\begin{itemize}
  \item A \textbf{query library}: pre-built Prolog assertion templates for common NF properties (RFC compliance, no port exhaustion, no silent drops)
  \item A \textbf{chain composition API}: standard interface for composing per-NF CFG-NC analyses into chain analyses
  \item A \textbf{result format}: standardized verification certificates that can be shared across NF deployments
\end{itemize}

\textbf{Strategy E --- Extract and cache behavioral summaries like Vigor's composable certificates}:\\
Vigor verifies library components once and reuses certificates. YP should maintain a \textbf{behavioral summary cache}: when a common eBPF helper (e.g., a specific Katran function) appears in a program, use the cached CFG-NC summary instead of re-analyzing the binary. This reduces analysis time for NFs built on common libraries.
\end{strategybox}

\begin{limitbox}
\begin{itemize}
  \item Vigor is \textbf{source-level}: developers must write NFs in Vigor's C framework. YP makes no such requirement.
  \item Vigor has \textbf{no chain verification}: it verifies individual NFs, not service chains. This is a critical gap that Yaksha-Prashna directly addresses.
  \item Vigor's KLEE-based approach scales poorly to \textbf{complex NFs} with large state spaces (path explosion).
  \item Vigor targets \textbf{DPDK}, not eBPF. The eBPF ecosystem (with its kernel verifier, map types, and XDP/TC hooks) requires a completely different verification approach.
\end{itemize}
\end{limitbox}

% ------------------------------------------------------------
\subsection{Connecting VigNAT and Vigor to Our Research Agenda}
% ------------------------------------------------------------

VigNAT and Vigor together establish the most important research insight for improving Yaksha-Prashna:

\begin{insightbox}
\textbf{The fundamental decomposition principle}: Separate NF verification into (1) \emph{stateful data structure correctness} (hard, requires formal methods) and (2) \emph{packet processing logic correctness} (easier, amenable to symbolic execution or static analysis). Prove each independently, then compose.

For YP, this translates to: separate \emph{BPF map state modeling} (hard --- needs formal map algebra) from \emph{packet field processing analysis} (easier --- YP's current CFG-NC is already good at this). Prove each part with the appropriate technique, then compose via YP's chain analysis framework.
\end{insightbox}

The other key lesson: Vigor proves that \emph{push-button usability} is achievable without sacrificing formal guarantees. YP is already push-button (just load the bytecode). The goal is to match Vigor's guarantee strength while maintaining YP's usability advantage.


% ============================================================
%  Section: BeePL, AppNet, eNetSTL Analysis
% ============================================================

\section{Research Analysis: BeePL, AppNet, and eNetSTL}
\label{sec:beepl-appnet-enetSTL}

% ------------------------------------------------------------
\subsection{BeePL: Correct-by-Construction Kernel Extensions}
% ------------------------------------------------------------

\subsection*{Paper: BeePL: Correct-by-Construction Kernel Extensions}
\textbf{Year}: 2025 \quad \textbf{arXiv}: 2502.01234 \quad \textbf{Authors}: Not listed (preprint)

\subsubsection{First-Pass Intuition}

The eBPF verifier is the gatekeeper of the Linux kernel. Every eBPF program must pass its safety checks before being loaded. But the verifier has a well-known problem: it is \emph{overly conservative}. Correct programs sometimes get rejected (false positives from the verifier's perspective); and rarely but seriously, buggy programs slip through because the verifier checks the wrong thing (unsound paths, as shown by Heimdall).

BeePL takes a radically different approach: \textbf{don't fix the verifier, eliminate the need for it}. BeePL is a new programming language (a DSL --- Domain-Specific Language) for writing eBPF programs in which \emph{unsafe constructs are not expressible in the grammar}. You cannot write an unbounded loop in BeePL. You cannot perform unchecked pointer arithmetic. You cannot call a helper without a null check. If your program compiles, it is safe --- by construction.

BeePL then compiles to eBPF bytecode using a \emph{type-preserving compiler}: the safety types from the source language are preserved through compilation, so the resulting bytecode trivially passes the eBPF verifier (and can be proven to do so).

\subsubsection{Classification as NF Validation}

BeePL is at the \emph{safety} end of the NF validation spectrum. It guarantees memory safety, pointer safety, bounded loops, and type-safe helper calls --- but it does \emph{not} guarantee functional correctness (that the NF correctly classifies, routes, or modifies packets according to a policy). A correct-by-construction BeePL firewall might be memory-safe but still implement the wrong firewall rules. This is a critical distinction.

\begin{insightbox}
BeePL gives you: \textbf{safety guarantees for free} (no runtime bugs, no verifier rejections). \\
YP gives you: \textbf{behavioral correctness assertions} (does the NF do the right thing?). \\
Both are needed. They are complementary, not competing.
\end{insightbox}

\subsubsection{Technical Mechanism}

\paragraph{The Type System.}
BeePL's type system has three key components:
\begin{enumerate}
  \item \textbf{Pointer types}: distinct types for packet pointers (\texttt{pkt\_ptr}), map value pointers (\texttt{map\_val\_ptr}), and stack pointers (\texttt{stack\_ptr}). Each has different dereference rules, enforced at compile time.
  \item \textbf{Option types}: BPF map lookups return \texttt{Option<Value>} (either \texttt{Some(v)} or \texttt{None}). You must pattern-match on the result before using it --- null pointer dereferences are impossible by type.
  \item \textbf{Loop typing}: Loops require a lexicographic termination measure (a proof that the loop terminates) expressed as a type annotation. Unbounded loops cannot be typed.
\end{enumerate}

\paragraph{Type-Preserving Compilation.}
The BeePL compiler translates source-level types to eBPF type annotations in the bytecode. The Linux eBPF verifier reads these annotations and immediately accepts the program --- no need for the verifier to re-derive safety properties from scratch.

\paragraph{Progress + Preservation.}
BeePL's type system is proven sound via \emph{progress and preservation} lemmas (standard type theory):
\begin{itemize}
  \item \textbf{Progress}: Well-typed programs don't get stuck.
  \item \textbf{Preservation}: Types are preserved through evaluation.
\end{itemize}
Together, these guarantee that a well-typed BeePL program never crashes at runtime.

\textbf{Tools}: Custom DSL + type checker, custom compiler (to eBPF bytecode), formal proof of type soundness (Coq or Isabelle).

\begin{analogybox}
\textbf{BeePL's type system $\leftrightarrow$ YP's verifier-derived type information}

BeePL assigns types \emph{before} compilation; the eBPF verifier propagates types \emph{after} compilation. YP could read the \emph{verifier's output} (BTF metadata, BPF type annotations embedded in ELF) to recover the same type information that BeePL starts with.

Specifically:
\begin{itemize}
  \item BeePL's \texttt{pkt\_ptr} $\leftrightarrow$ eBPF verifier's \texttt{PTR\_TO\_PACKET} register type $\leftrightarrow$ YP's packet field access tracking
  \item BeePL's \texttt{Option<Value>} for map lookups $\leftrightarrow$ eBPF verifier's \texttt{PTR\_TO\_MAP\_VALUE\_OR\_NULL} $\leftrightarrow$ YP's map read operation nodes
  \item BeePL's termination measure $\leftrightarrow$ eBPF verifier's loop bound annotations (in BTF ext) $\leftrightarrow$ YP's loop detection in CFG-NC
\end{itemize}
YP can mine BTF metadata to enrich CFG-NC with type information that BeePL explicitly tracks. This expands the set of analyzable fields beyond the current 6.
\end{analogybox}

\begin{strategybox}
\textbf{Strategy F --- BTF-driven CFG-NC enrichment}:\\
eBPF programs compiled with debug info contain BTF (BPF Type Format) metadata: register types at each instruction, map key/value type information, function signatures. YP currently ignores this metadata. Extracting BTF types would:
\begin{enumerate}
  \item Automatically identify which registers hold packet pointers vs. map value pointers vs. stack data
  \item Expand tracked field coverage beyond the current 6 fields to all typed packet fields
  \item Enable type-aware Prolog queries (e.g., ``does this program ever store a packet pointer in a map?'' --- a potential data exfiltration query)
\end{enumerate}

\textbf{Strategy G --- Treat BeePL programs as trusted inputs}:\\
If an NF was written in BeePL, YP can skip its safety checks (guaranteed by BeePL's type system) and focus exclusively on functional correctness queries. This is a compositional trust model: BeePL proves safety; YP proves correctness.
\end{strategybox}

\begin{limitbox}
\begin{itemize}
  \item BeePL requires developers to \textbf{rewrite existing NFs in a new language}. No migration path for the thousands of existing eBPF NFs.
  \item BeePL \textbf{does not prove functional correctness}: a BeePL firewall may be memory-safe but enforce the wrong policy.
  \item \textbf{No chain verification}: BeePL is single-program only. Chain-level properties (ordering, stateful carry-over) are outside its scope.
  \item BeePL is a \textbf{research prototype}, not production-ready. Adoption is uncertain.
\end{itemize}
\end{limitbox}

% ------------------------------------------------------------
\subsection{AppNet: Semantic-Aware Optimization of Application Network Functions}
% ------------------------------------------------------------

\subsection*{Paper: AppNet: High-Level Programming for Application Networks}
\textbf{Year}: 2025 \quad \textbf{Venue}: USENIX NSDI 2025 \quad \textbf{Authors}: Xiangfeng Zhu et al., University of Washington / Duke / UCLA

\subsubsection{First-Pass Intuition}

Modern microservice applications deploy a layer of software ``middleboxes'' between every pair of services: authentication filters, rate limiters, circuit breakers, load balancers, tracing sidecars. These are collectively called \emph{Application Network Functions} (ANFs). They are typically implemented as Envoy or Istio filters.

The problem: each ANF adds latency (20--200 ms each), and naive placement and ordering creates unnecessary overhead. But you can't just reorder or merge ANFs arbitrarily --- changing the order of ``authentication before rate limiting'' vs.\ ``rate limiting before authentication'' changes which requests are processed and has security implications.

AppNet's idea: use \textbf{Z3 SMT equivalence checking} to automatically prove whether a proposed ANF chain optimization (reordering, merging, bypassing) is semantically safe. If Z3 says the optimized chain is equivalent to the original for all possible requests, the optimization is provably safe. Result: up to 82\% latency reduction, 75\% CPU savings.

\subsubsection{Classification as NF Validation}

AppNet is primarily an \emph{NF chain optimization} paper, but its core technical contribution is \textbf{formal NF chain equivalence checking} --- which is directly NF validation at the chain level. The equivalence checking machinery (symbolic state transformers + Z3) is exactly what Yaksha-Prashna needs for its chain verification module.

\subsubsection{Technical Mechanism}

\paragraph{Symbolic State Transformers.}
Each ANF is modeled as a \emph{symbolic state transformer}: a mathematical function $f: (\text{Request} \times \text{State}) \to (\text{Response} \times \text{State}')$. This models how the ANF transforms both the request and the shared state on each invocation.

\paragraph{Chain Composition.}
A chain of ANFs is modeled as the sequential composition of their state transformers: $f_n \circ \cdots \circ f_2 \circ f_1$. The composed function maps any input request and initial state to the final response and final state.

\paragraph{Equivalence Checking via Z3.}
Two chain configurations (original $C_1$, optimized $C_2$) are equivalent iff:
\[ \forall r \in \text{Requests},\ s \in \text{States}: C_1(r, s) = C_2(r, s) \]
This is encoded as a Z3 formula. Z3 either proves UNSAT (no counterexample exists --- chains are equivalent) or produces a concrete (request, state) pair that demonstrates the behavioral difference.

\paragraph{Cost-Based Optimization.}
Given proven-equivalent configurations, AppNet selects the cheapest one based on a cost model: in-library execution $<$ sidecar $<$ middlebox. Backends include gRPC interceptors, Envoy Wasm filters, and eBPF programs.

\textbf{Tools}: Z3 SMT solver, AppNet DSL (match-action on RPC fields), Envoy/Istio configuration backends.

\begin{analogybox}
\textbf{AppNet's symbolic state transformer composition $\leftrightarrow$ YP's CFG-NC chain composition}

AppNet and YP share the same problem at different layers:
\begin{itemize}
  \item AppNet reasons about L7 (RPC/HTTP) ANF chains using symbolic state transformers
  \item YP reasons about L3/L4 (eBPF XDP/TC) NF chains using CFG-NC composition
\end{itemize}

The mathematical structure is identical: both compose ``NF models'' (state transformers / CFG-NCs) sequentially and check properties of the composed result. AppNet's Z3-based equivalence checking is the formal foundation YP needs for chain-level correctness queries.

What we find particularly valuable: AppNet's \textbf{commutativity checking} (can NF$_i$ and NF$_j$ be swapped?) is directly implementable in YP's Prolog engine as: ``do the two orderings produce the same packet decisions and map states for all inputs?''
\end{analogybox}

\begin{strategybox}
\textbf{Strategy H --- Adopt AppNet's state transformer model for YP's chain composition}:\\
Formalize each eBPF NF in YP's chain as a symbolic state transformer $f_i: (\text{BPF maps}, \text{packet}) \to (\text{BPF maps}', \text{packet}', \text{XDP action})$. Chain composition becomes sequential function composition, enabling:
\begin{enumerate}
  \item \textbf{Chain equivalence queries}: ``Are these two NF orderings equivalent?''
  \item \textbf{Commutativity detection}: ``Can NF$_i$ and NF$_j$ be safely reordered?''
  \item \textbf{Dead NF detection}: ``Does NF$_k$ ever change the packet/state in this chain?''
\end{enumerate}

\textbf{Strategy I --- Use AppNet's Z3 backend for complex chain queries}:\\
When YP's Prolog engine cannot resolve a complex chain equivalence query symbolically, fall through to Z3 with the AppNet-style encoding. This gives YP a ``heavy solver'' fallback for queries that require full symbolic reasoning.
\end{strategybox}

\begin{limitbox}
\begin{itemize}
  \item AppNet targets \textbf{L7 RPC/HTTP ANFs} (Envoy, Istio), not \textbf{L3/L4 eBPF kernel programs}. The state spaces are fundamentally different.
  \item AppNet requires NFs to be written in \textbf{AppNet's DSL}. YP needs no such restriction.
  \item AppNet proves \textbf{equivalence between configurations} (both correct), not \textbf{correctness against a specification}. YP needs the latter.
  \item Z3 may \textbf{time out} on complex ANF compositions with large state spaces.
\end{itemize}
\end{limitbox}

% ------------------------------------------------------------
\subsection{eNetSTL: Safe and Efficient In-Kernel eBPF Network Function Library}
% ------------------------------------------------------------

\subsection*{Paper: eNetSTL: Towards an In-Kernel Library for High-Performance eBPF-based Network Functions}
\textbf{Year}: 2025 \quad \textbf{Venue}: EuroSys 2025 \quad \textbf{Authors}: Bin Yang, Dian Shen, Junxue Zhang et al.

\subsubsection{First-Pass Intuition}

eBPF programs are powerful but limited: they cannot call arbitrary kernel functions, use complex data structures, or invoke SIMD operations. This forces developers to re-implement common NF primitives (SIMD hashing, parallel key comparison, counting sketches) in restricted eBPF C, leading to slow, redundant, and hard-to-verify code.

eNetSTL is the solution: a \emph{C++ standard template library for eBPF programs}, implemented as a Rust kernel module. It provides SIMD-accelerated algorithms (bitwise/POPCNT, parallel hashing, parallel key comparison) and data structures (list buckets, random number pools) as \emph{kernel functions} (\texttt{kfunc}) that eBPF programs can call.

The key innovation for verification: eNetSTL introduces \textbf{metadata-assisted eBPF verification}. The library's \texttt{kfunc} declarations carry special annotations (\texttt{\_\_sz}, \texttt{\_\_nullable}, \texttt{\_\_uninit}) that guide the eBPF verifier to accept safe programs that would otherwise be rejected. The verifier reads these annotations from BTF metadata and applies them directly to type checking.

\subsubsection{Technical Mechanism}

\paragraph{kfunc + BTF Annotations.}
eNetSTL functions are declared in Rust with annotations:
\begin{itemize}
  \item \texttt{\_\_sz}: ``this parameter specifies the size of the associated buffer parameter'' --- enables safe buffer access without verifier guessing
  \item \texttt{\_\_nullable}: ``this pointer may be NULL, check before use'' --- verifier enforces null check
  \item \texttt{\_\_uninit}: ``this buffer is uninitialized, must be written before read'' --- prevents uninitialized memory reads
\end{itemize}
The Rust BTF type graph propagates these from the module to the verifier at load time.

\paragraph{Library Taxonomy.}
eNetSTL covers 7 NF operation categories:
key-value query, membership testing, packet classification, load balancing, counting, sketching, and queuing. Each is provided as a verified, SIMD-optimized kernel function.

\paragraph{Performance.}
eNetSTL achieves 14.6--75.4\% throughput improvement over pure eBPF equivalents, within 3.42\% average of native kernel implementations.

\textbf{Integrated NFs}: Katran (Facebook load balancer), PolyCube, RakeLimit, Count-Min Sketch, Nitro Sketch.

\begin{analogybox}
\textbf{eNetSTL's behavioral contracts $\leftrightarrow$ YP's inter-procedural analysis}

When YP analyzes an eBPF program that calls a \texttt{kfunc} from eNetSTL, it currently has two choices: (1) treat the call as opaque (lose behavioral information), or (2) try to analyze the kernel module implementation (expensive, requires source).

eNetSTL's BTF annotations provide a \textbf{third option}: read the library's declared behavioral contracts directly from BTF metadata. YP can trust eNetSTL's annotations (null-safety, size contracts, etc.) as pre-verified behavioral summaries, enabling inter-procedural analysis of \texttt{kfunc} calls without analyzing the kernel module.

This is directly analogous to Vigor's ``trusted library'' approach: verify the library once (eNetSTL does this), reuse the contracts everywhere (YP does this by reading BTF).
\end{analogybox}

\begin{strategybox}
\textbf{Strategy J --- YP's Library Contract Database}:\\
Build a curated database of behavioral contracts for common eBPF library functions and kfuncs:
\begin{enumerate}
  \item \textbf{eNetSTL kfuncs}: behavioral contracts directly from BTF annotations
  \item \textbf{Common helpers}: \texttt{bpf\_map\_lookup\_elem}, \texttt{bpf\_redirect}, \texttt{bpf\_xdp\_adjust\_head}, etc. --- pre-computed formal models
  \item \textbf{Katran/Cilium functions}: behavioral summaries for widely-used eBPF components
\end{enumerate}
When YP encounters a known function, substitute the pre-computed contract instead of re-analyzing. This eliminates the ``library function opacity'' problem and significantly reduces analysis time.

\textbf{Strategy K --- Treat eNetSTL's 7-category taxonomy as YP's NF classification scheme}:\\
eNetSTL's 7 categories (key-value, membership, classification, load balancing, counting, sketching, queuing) can serve as YP's default NF taxonomy for organizing Prolog assertion templates and evaluation benchmarks.
\end{strategybox}

\begin{limitbox}
\begin{itemize}
  \item eNetSTL focuses on \textbf{API-boundary safety}, not functional correctness of NF policy logic.
  \item Library contracts are \textbf{trusted by assumption} --- if eNetSTL has a bug, YP's contract-based analysis is wrong.
  \item \textbf{No chain verification}: eNetSTL is a library, not a chain analysis framework.
  \item Adoption requires developers to \textbf{refactor existing NFs} to use eNetSTL APIs.
\end{itemize}
\end{limitbox}

% ------------------------------------------------------------
\subsection{Synthesis: Three Papers, One Message}
% ------------------------------------------------------------

BeePL, AppNet, and eNetSTL each attack a different layer of the NF verification challenge:

\begin{center}
\begin{tabular}{llll}
\toprule
\textbf{Paper} & \textbf{Layer} & \textbf{What it guarantees} & \textbf{What it misses} \\
\midrule
BeePL & Source (DSL) & Memory/pointer safety & Functional correctness \\
AppNet & L7 chain & Chain equivalence & L3/L4, spec compliance \\
eNetSTL & Library API & API safety contracts & NF policy correctness \\
\midrule
\textbf{YP} & \textbf{L3/L4 bytecode} & \textbf{Behavioral assertions} & \textbf{(all three above, to varying degrees)} \\
\bottomrule
\end{tabular}
\end{center}

\begin{insightbox}
What these three papers collectively teach us: \textbf{NF verification is a stack problem}. Safety (BeePL), library correctness (eNetSTL), chain equivalence (AppNet), and behavioral policy correctness (YP) are different layers of the same problem. The research agenda for YP should focus on the policy correctness layer while incorporating contract-based trust from the layers below.
\end{insightbox}


% ============================================================
%  Section: Demystifying eBPF + FlowMage Analysis
% ============================================================

\section{Research Analysis: Demystifying eBPF Performance and FlowMage}
\label{sec:demystifying-flowmage}

% ------------------------------------------------------------
\subsection{Demystifying Performance of eBPF Network Applications}
% ------------------------------------------------------------

\subsection*{Paper: Demystifying Performance of eBPF Network Applications}
\textbf{Year}: 2025 \quad \textbf{Venue}: ACM CoNEXT 2025 (PACMNET Vol.\ 3, Article 16)\\
\textbf{DOI}: \url{https://dl.acm.org/doi/10.1145/3700898}\\
\textbf{Authors}: Farbod Shahinfar, Sebastiano Miano, Aurojit Panda, Gianni Antichi\\
\textbf{Artifact}: \url{https://github.com/bpf-endeavor/bpf-app-offload-measurement}

\subsubsection{First-Pass Intuition}

There is a widespread assumption in the eBPF community: moving packet processing into the kernel (using XDP, TC, or SK\_SKB hooks) always makes things faster. Cilium, Katran, Cloudflare's firewall, and BMC (BPF Memcached cache) all leverage eBPF on this assumption. This paper asks: \emph{is it actually true?}

The answer is: \textbf{it depends, and often the answer is no}. eBPF offloading delivers performance gains only under specific conditions. When those conditions are not met, eBPF can actively \emph{hurt} performance --- both for the offloaded application and for co-located workloads.

The paper builds a precise analytical model (the $\omega/\rho$ model) that predicts exactly when eBPF offloading is beneficial:

\begin{itemize}
  \item $\omega$: the fraction of requests fully handled by the eBPF fast path (never reaching userspace)
  \item $\rho$: the overhead ratio of the eBPF program (cycles with program / cycles without)
  \item \textbf{Offloading is net-positive only when}: $\omega > \rho - 1$
\end{itemize}

If too few requests can be served in-kernel (low $\omega$), or the eBPF program is expensive (high $\rho$), offloading degrades throughput. This model is validated against real measurements of BMC and AF\_XDP-based systems on 100 Gbps hardware.

\subsubsection{Classification as NF Validation}

This paper is \emph{not} a correctness verification paper --- it does not check whether NFs forward packets correctly. However, it is deeply relevant to NF validation for three reasons:

\begin{enumerate}
  \item It establishes \textbf{performance ground truth} for eBPF NFs --- a prerequisite for performance-level behavioral assertions in YP.
  \item It identifies \textbf{verifier-induced overhead} (loop bound-checks adding 5--15\% throughput penalty) --- a correctness-adjacent concern that YP can detect statically.
  \item It documents \textbf{performance isolation violations} --- when one eBPF NF degrades another's performance, which is a behavioral property YP could assert about.
\end{enumerate}

\subsubsection{Technical Mechanism}

\paragraph{The $\omega/\rho$ Model.}
The performance model is derived from first principles. Let $T_0$ be baseline throughput (no eBPF program). With an eBPF program:
\begin{align}
  T_{\text{eBPF}} = \frac{T_0}{\rho - \omega}
\end{align}
eBPF offloading is net-positive when $T_{\text{eBPF}} > T_0$, i.e., $\omega > \rho - 1$. The model is validated with $<$10\% error on real workloads.

\paragraph{Benchmarking Methodology.}
The paper conducts a systematic 10-step benchmarking campaign:
\begin{enumerate}
  \item Hook overhead: minimal pass-through eBPF program at XDP, TC, SK\_SKB --- measures irreducible entry cost
  \item Time-to-hook: latency from NIC arrival to hook invocation at each layer
  \item Round-trip latency: echo server at each hook (XDP vs.\ TC vs.\ SK\_SKB vs.\ userspace)
  \item Map API overhead: Array vs.\ Hash vs.\ Ring Buffer access latency at varying sizes
  \item Compilation overhead: Clang/LLVM JIT code quality vs.\ native
  \item Verifier overhead: loop bound-checks in tight loops (5--15\% throughput impact)
  \item System interference: co-located eBPF NF degrading unrelated workloads
  \item Packet trimming: \texttt{bpf\_xdp\_adjust\_tail} reducing DMA footprint
  \item BMC case study: eBPF cache for Memcached at XDP hook
  \item xsk\_cache: AF\_XDP-based key-value store
\end{enumerate}

\paragraph{Key Findings.}
\begin{itemize}
  \item XDP native hook: $\sim$300 ns overhead; TC: $\sim$500 ns; SK\_SKB: $>$1 $\mu$s
  \item eBPF Hash map lookup: 100--200 ns; Array map: 20--50 ns
  \item JIT-compiled eBPF bulk operations are \textbf{2--5$\times$ slower} than equivalent native code (no SIMD support in eBPF JIT)
  \item Verifier-inserted loop bound-checks cause \textbf{5--15\% throughput loss}
  \item BMC achieves net-positive offloading only at $>$60\% cache hit rate ($\omega > 0.6$)
\end{itemize}

\textbf{Testbed}: Two servers, 100 Gbps Mellanox ConnectX-6 NICs, Intel Xeon Silver 4310, Linux 6.8.0-rc7, Clang 14 with eBPF ISA v3 ($-$mcpu=v3).

\begin{analogybox}
\textbf{The $\omega/\rho$ model $\leftrightarrow$ YP's static hook classification}

YP's CFG-NC already records which hook type each eBPF program uses (XDP, TC, SK\_SKB). This paper provides the quantitative meaning of that classification:

\begin{itemize}
  \item XDP programs have $\rho \approx 1.1$--$1.3$ (low overhead) --- need only $\omega > 0.1$--$0.3$ to be net-positive
  \item TC programs have $\rho \approx 1.3$--$1.6$ --- need higher hit rates
  \item SK\_SKB programs have $\rho > 2$ in many cases --- rarely net-positive for latency-sensitive NFs
\end{itemize}

YP can \emph{statically estimate $\rho$} from the bytecode: count map lookups (100--200 ns each), helper calls, conditional branches, and match against the hook overhead. This gives a static performance bound without running the NF.

Moreover: YP already detects loop bounds from bytecode (the verifier requires them). The paper shows these bounds incur 5--15\% overhead. YP can query: ``does this loop have a suboptimal bound annotation that degrades throughput?''
\end{analogybox}

\begin{strategybox}
\textbf{Strategy L --- Static performance assertion layer for YP}:\\
Integrate the $\omega/\rho$ model into YP's Prolog query vocabulary:
\begin{lstlisting}[language=Prolog]
% New YP Prolog query:
performance_viable(Program, Workload) :-
    hook_type(Program, Hook),
    hook_overhead(Hook, Rho),
    expected_hit_rate(Workload, Omega),
    Omega > Rho - 1.
\end{lstlisting}
At static analysis time, YP computes $\rho$ from the bytecode (map lookup count $\times$ cost model + branch cost) and checks the performance viability condition against user-specified workload parameters.

\textbf{Strategy M --- Verifier loop overhead detection}:\\
YP reads loop bound annotations from BTF metadata. When a loop has a bound much larger than the actual worst-case flow count, YP flags this as a ``verifier overhead inefficiency'' --- the developer specified a conservative bound that causes unnecessary overhead. Suggest a tighter bound.

\textbf{Strategy N --- NF co-location interference queries}:\\
When analyzing an NF chain, YP can check for shared per-CPU resources (same CPU affinity, same map access patterns) and flag potential co-location interference. Query: ``do NF$_i$ and NF$_j$ in this chain compete for the same per-CPU map?''
\end{strategybox}

\begin{limitbox}
\begin{itemize}
  \item This paper provides \textbf{no functional correctness} analysis --- only performance.
  \item The $\omega/\rho$ model assumes \textbf{stationary workloads}. Real traffic has bursts that invalidate static performance predictions.
  \item Measurements are \textbf{hardware-specific} (ConnectX-6 NICs, Intel Xeon Silver) --- performance bounds may differ significantly on different hardware.
  \item \textbf{No NF chain analysis}: all experiments are single-NF, single-hook.
\end{itemize}
\end{limitbox}

% ------------------------------------------------------------
\subsection{FlowMage: LLM-Based Static Analysis for Stateful NF Chain Optimization}
% ------------------------------------------------------------

\subsection*{Paper: FlowMage: LLM-Based Static Analysis for Stateful NF Chain Optimization}
\textbf{Year}: 2024 \quad \textbf{Venue}: EuroMLSys (ACM Workshop on ML for Systems) \quad \textbf{Authors}: Hamid Ghasemi et al.

\subsubsection{First-Pass Intuition}

Deploying a stateful NF chain on a multi-core server requires careful configuration of \textbf{RSS (Receive Side Scaling)} --- the mechanism that distributes incoming network flows across CPU cores. The problem: stateful NFs require that all packets of the same flow be processed on the same CPU core (so the NF's per-flow state machine sees all packets). If RSS is misconfigured, flow packets arrive on different cores, the state machine gets confused, and the NF fails silently.

Manually configuring RSS for complex NF chains is error-prone. FlowMage's solution: use \textbf{GPT-4 to read NF source code} and automatically extract which packet header fields define a ``flow'' for each NF in the chain. From these extracted \emph{affinity constraints}, compute the minimum RSS hash key that ensures flow coherence for the entire chain.

This is a direct application of LLMs to NF semantic property extraction --- and it works with surprisingly high accuracy.

\subsubsection{Classification as NF Validation}

FlowMage validates \textbf{configuration correctness}: given an RSS configuration, does it satisfy the affinity requirements of the NF chain? This is NF validation at the deployment configuration level, not the packet processing logic level. It answers: ``will the infrastructure correctly route flows to this NF chain?'' rather than ``does the NF correctly process flows?''

\begin{insightbox}
\textbf{FlowMage is to RSS configuration what YP is to NF behavior.} FlowMage validates the \emph{routing} is correct; YP validates the \emph{processing} is correct. Both are needed for end-to-end NF correctness.
\end{insightbox}

\subsubsection{Technical Mechanism}

\paragraph{Step 1 --- LLM Property Extraction.}
For each NF in the chain (FastClick element or VPP plugin written in C++), send a structured prompt to GPT-4:
\begin{itemize}
  \item ``Which packet header fields does this NF read?''
  \item ``Which fields define a flow for this NF's state machine?''
  \item ``Would processing two packets of the same flow on different cores cause errors?''
\end{itemize}
GPT-4 returns structured JSON: \texttt{\{reads: [src\_ip, dst\_ip, src\_port, ...], affinity: [src\_ip, dst\_ip, src\_port, dst\_port, proto], is\_stateful: true\}}

\paragraph{Step 2 --- Affinity Constraint Computation.}
Compute the chain-wide affinity requirement: $\text{ChainAffinity} = \bigcup_{i} \text{Affinity}_i$. This is the minimum RSS hash key that satisfies all NFs simultaneously.

\paragraph{Step 3 --- RSS Validation.}
Compare the proposed RSS configuration's hash key against $\text{ChainAffinity}$. If the hash key is not a superset of $\text{ChainAffinity}$, the configuration is invalid --- flag it and suggest the correct hash key.

\textbf{Tools}: GPT-4 API, FastClick (C++ Click framework), VPP (Vector Packet Processor), structured JSON prompting.

\begin{analogybox}
\textbf{FlowMage using GPT-4 $\leftrightarrow$ YP using static dataflow analysis}

This is the most direct analogy in our entire literature survey:

\begin{center}
\begin{tabular}{lll}
\toprule
\textbf{Property} & \textbf{FlowMage (LLM)} & \textbf{Yaksha-Prashna (static)} \\
\midrule
What fields does NF read? & GPT-4 reads source code & CFG-NC tracks packet field accesses \\
What defines a flow? & GPT-4 infers affinity & Prolog query on CFG-NC reads \\
Is the config valid? & Affinity set intersection & Prolog assertion check \\
Guarantee & Probabilistic (LLM) & Deterministic (dataflow) \\
Requires source code? & Yes & \textbf{No} \\
\bottomrule
\end{tabular}
\end{center}

\textbf{YP already does, deterministically and formally, what FlowMage does probabilistically with LLMs.} This is a strong positioning argument for YP in the literature.

The key insight: FlowMage's affinity constraint formulation (which fields define a ``flow'') maps exactly to YP's CFG-NC packet field access tracking. YP should add a \textbf{flow affinity query} to its Prolog vocabulary.
\end{analogybox}

\begin{strategybox}
\textbf{Strategy O --- Add flow affinity queries to YP's Prolog engine}:\\
Implement the following Prolog query in YP:
\begin{lstlisting}[language=Prolog]
% Flow affinity: fields that define a unique flow for this NF
flow_affinity_fields(Program, Fields) :-
    findall(F, (stateful_read(Program, F), 
                map_key_component(Program, F)), Fields).

% Chain affinity: union of all NF affinity fields in chain  
chain_affinity(Chain, ChainAffinity) :-
    maplist(flow_affinity_fields, Chain, PerNFAffinity),
    foldl(union, PerNFAffinity, [], ChainAffinity).

% RSS validation: does the hash key cover chain affinity?
rss_valid(Chain, RSSHashKey) :-
    chain_affinity(Chain, Required),
    subset(Required, RSSHashKey).
\end{lstlisting}
This enables YP to \textbf{formally validate RSS configurations} --- strictly superior to FlowMage's LLM approach.

\textbf{Strategy P --- Use FlowMage as a positioning benchmark}:\\
FlowMage's accuracy on FastClick elements (high GPT-4 accuracy, low false negatives) provides a comparison baseline. Demonstrate that YP computes flow affinity with \emph{100\% accuracy} (formal dataflow) vs.\ FlowMage's probabilistic LLM accuracy. This is a concrete experimental contribution.
\end{strategybox}

\begin{limitbox}
\begin{itemize}
  \item FlowMage requires \textbf{C++ source code} (FastClick/VPP). YP analyzes compiled eBPF bytecode.
  \item FlowMage validates \textbf{configuration correctness only}, not packet processing logic.
  \item LLM extraction is \textbf{probabilistic}: GPT-4 can misclassify affinity fields, causing false negatives in validation.
  \item FlowMage is specific to \textbf{FastClick/VPP frameworks}, not applicable to eBPF.
  \item \textbf{Workshop paper} (EuroMLSys) --- limited peer review depth.
\end{itemize}
\end{limitbox}

% ------------------------------------------------------------
\subsection{Synthesis: Performance + Configuration as Validation Dimensions}
% ------------------------------------------------------------

These two papers expand our understanding of what ``NF validation'' means:

\begin{insightbox}
NF validation is not just about functional correctness (does the firewall correctly implement the ACL rules?). It also encompasses:
\begin{itemize}
  \item \textbf{Performance correctness} (Demystifying eBPF): does the NF meet its performance SLA? Will it degrade co-located workloads?
  \item \textbf{Configuration correctness} (FlowMage): is the NF's deployment infrastructure configured to route flows correctly to the NF chain?
\end{itemize}
YP currently addresses only functional correctness. The research agenda should explore adding \emph{performance assertions} and \emph{configuration validation} as first-class query types in YP's Prolog engine.
\end{insightbox}


% ============================================================
%  Section: NetSMC + Full-Lifecycle Intent Analysis
% ============================================================

\section{Research Analysis: NetSMC and Full-Lifecycle Intent-Driven Verification}
\label{sec:netSMC-fulllifecycle}

% ------------------------------------------------------------
\subsection{NetSMC: A Custom Symbolic Model Checker for Stateful Network Verification}
% ------------------------------------------------------------

\subsection*{Paper: NetSMC: A Custom Symbolic Model Checker for Stateful Network Verification}
\textbf{Year}: 2020 \quad \textbf{Venue}: USENIX NSDI 2020\\
\textbf{DOI}: \url{https://www.usenix.org/conference/nsdi20/presentation/yuan}\\
\textbf{Authors}: Yifei Yuan, Soo-Jin Moon, Sahil Uppal, Limin Jia, Vyas Sekar (CMU)

\subsubsection{First-Pass Intuition}

Verifying a single NF is hard. Verifying a \emph{chain} of stateful NFs is exponentially harder. A firewall's state depends on which connections it has seen; a load balancer's state depends on which backends are healthy; an IDS's state depends on which signatures have matched. When these NFs are chained, their states interact: a packet that changes the firewall's state might change which backend the load balancer selects, which changes what the IDS sees.

Prior work (VMN, Kinetic) attempts to verify NF chains using general-purpose model checkers (NuSMV, PRISM). But these tools have no knowledge of network semantics, leading to enormous state spaces and verification times measured in hours or failing to terminate.

NetSMC's insight: \textbf{build a custom symbolic model checker that understands NF chain semantics}. By encoding domain-specific optimizations (containment checking, shared state factoring, per-NF state bound exploitation), NetSMC achieves \textbf{28--200$\times$ speedup} over VMN while verifying the same properties.

\subsubsection{Classification as NF Validation}

\textbf{NetSMC is the most directly relevant paper to Yaksha-Prashna's chain verification goals.} It solves exactly the problem YP faces: how to verify properties of stateful NF chains without exploding in state space. The difference is the input level: NetSMC takes formal NF models; YP takes compiled eBPF bytecode. NetSMC's algorithms can inform YP's chain analysis engine directly.

\subsubsection{Technical Mechanism}

\paragraph{Property Language: LTL over Packet Traces.}
NetSMC verifies Linear Temporal Logic (LTL) properties over packet traces. Examples:
\begin{itemize}
  \item $\mathbf{G}(\text{Admitted}(p) \Rightarrow \text{FW}(p) \wedge \text{LB}(p))$: ``Every admitted packet was processed by the firewall and load balancer''
  \item $\mathbf{G}(\text{External}(p) \Rightarrow \neg\text{Internal}(p)\ \mathbf{U}\ \text{Authenticated}(p))$: ``External packets don't reach internal resources until authenticated''
  \item $\mathbf{G}(\text{Drop}(p) \Rightarrow \text{FW\_Reject}(p) \vee \text{IDS\_Reject}(p))$: ``Dropped packets are dropped by FW or IDS, not silently lost''
\end{itemize}
These are temporal properties over \emph{sequences} of packet processing events --- something YP's current Prolog engine cannot express.

\paragraph{The Containment Optimization.}
NetSMC's primary speedup comes from \emph{containment checking}: rather than exploring all reachable states of the NF chain, NetSMC checks whether the set of states reachable under the negation of the property ($\neg \phi$) is contained in the set of states where the property is already known to hold ($\phi$). If it is, the property holds without full state space exploration.

Formally: let $R$ be reachable states and $G$ be the set where $\phi$ holds. If $R \subseteq G$, the property holds. NetSMC computes this containment check symbolically (using BDDs + Z3) rather than explicitly enumerating states.

\paragraph{NF State Model.}
Each NF is modeled as a finite-state transducer: a state machine with:
\begin{itemize}
  \item \textbf{State variables}: e.g., connection table (for FW), backend health (for LB), alert counter (for IDS)
  \item \textbf{Input}: packet headers
  \item \textbf{Transition function}: how state evolves on each packet
  \item \textbf{Output}: packet decision (forward/drop) + modified headers
\end{itemize}
The chain model is the \emph{synchronized product} of individual NF state machines.

\paragraph{Shared State Handling.}
NetSMC explicitly models shared state between NFs (e.g., a shared connection table consulted by both a firewall and an IDS). Shared state variables are factored out and maintained once, rather than duplicated in each NF's state machine.

\textbf{Tools}: Custom model checker (NSMC), Z3, BDDs (Binary Decision Diagrams), LTL model theory.\\
\textbf{NFs verified}: Stateful firewall, TCP state-tracking IDS, load balancer (in chains).\\
\textbf{Performance}: 28--200$\times$ faster than VMN on the same properties.

\begin{analogybox}
\textbf{NetSMC's LTL over packet traces $\leftrightarrow$ YP's Prolog behavioral assertions}

YP's Prolog assertions answer single-packet questions:
\begin{itemize}
  \item ``Does this NF ever drop packets with \texttt{src\_port=80}?'' (Prolog: possible)
  \item ``Does NF$_1$'s output reach NF$_2$ with field $F$ unmodified?'' (Prolog: possible)
\end{itemize}

NetSMC's LTL properties answer \emph{temporal} multi-packet questions:
\begin{itemize}
  \item ``If this packet is dropped, \emph{eventually} the connection will be retried and \emph{then} admitted''
  \item ``All packets from external sources must be authenticated \emph{before} reaching internal NFs''
  \item ``The load balancer's backend selection remains consistent \emph{until} a health-check event''
\end{itemize}

This is the \textbf{biggest capability gap in YP}: no temporal reasoning. NetSMC fills this gap for formal NF models; the research challenge is lifting NetSMC's LTL to the eBPF bytecode level.

\textbf{NetSMC's containment optimization $\leftrightarrow$ YP's path pruning}:
YP's static analysis already prunes infeasible paths. NetSMC's containment check is a more powerful version of this: prune states where the property is already known to hold. Adapting this to YP's Prolog engine would significantly reduce the query resolution cost for chain properties.
\end{analogybox}

\begin{strategybox}
\textbf{Strategy Q --- Extend YP's Prolog with LTL temporal meta-predicates}:\\
Implement temporal operators as Prolog meta-predicates over CFG-NC traces:
\begin{lstlisting}[language=Prolog]
% G(P): P holds for all packets in all execution traces
always(Program, P) :-
    forall(path(Program, Path), holds(P, Path)).

% F(P): P holds for some packet in some execution trace  
eventually(Program, P) :-
    path(Program, Path), holds(P, Path).

% P until Q: P holds until Q becomes true
until(Program, P, Q) :-
    path(Program, Path),
    split_at_first(Q, Path, Before, _After),
    forall(member(State, Before), holds(P, State)).

% Chain ordering: NF_i processes packet before NF_j
processes_before(Chain, NF_i, NF_j) :-
    chain_order(Chain, Order),
    before(NF_i, NF_j, Order).
\end{lstlisting}

\textbf{Strategy R --- Adopt NetSMC's containment optimization for YP's chain queries}:\\
For complex chain properties, implement symbolic containment checking:
\begin{enumerate}
  \item Compute the set of CFG-NC states where the property $\phi$ holds (symbolic)
  \item Compute the set of reachable states in the chain (symbolic)
  \item Check containment: if reachable $\subseteq$ $\phi$-satisfying, property holds without full enumeration
\end{enumerate}
This enables YP to scale to longer chains without state space explosion.

\textbf{Strategy S --- Use NetSMC's NF taxonomy as YP's chain evaluation benchmark}:\\
NetSMC's verified chains (FW+LB, FW+IDS, FW+LB+IDS) are the natural benchmark suite for YP's chain verification. Implement the same LTL properties as Prolog queries and compare verification times and guarantees.
\end{strategybox}

\begin{limitbox}
\begin{itemize}
  \item NetSMC requires \textbf{manually written formal NF models} (finite-state transducers). YP extracts models automatically from bytecode.
  \item NetSMC's formal models are \textbf{abstractions}: they may not faithfully capture all behaviors of real NF implementations (the abstraction gap).
  \item NetSMC is \textbf{not eBPF-aware}: it has no notion of BPF maps, XDP hooks, or eBPF helper calls.
  \item \textbf{State explosion} remains for very long chains or NFs with unbounded state (e.g., connection tables with no upper bound).
  \item NetSMC verifies \textbf{DPDK-era NFs} (stateful FW, LB, IDS as C programs), not the modern eBPF ecosystem.
\end{itemize}
\end{limitbox}

% ------------------------------------------------------------
\subsection{Full-Lifecycle Intent-Driven Network Verification}
% ------------------------------------------------------------

\subsection*{Paper: Full-Lifecycle Intent-Driven Network Verification}
\textbf{Year}: 2022 \quad \textbf{arXiv}: 2209.00999

\subsubsection{First-Pass Intuition}

Network operators manage complex NF deployments not by writing formal specifications, but by expressing \emph{intents}: high-level goals like ``all traffic from VLAN 10 must be inspected by the IDS before reaching the internet'' or ``ensure 99.9\% of traffic to service X is load-balanced across at least 3 backends.'' These intents are natural language (or semi-formal), not mathematical specifications.

The Full-Lifecycle Intent paper asks: \textbf{how do we ensure that a deployed NF chain actually implements an operator's intent, throughout the entire deployment lifecycle?}

The answer is a 4-phase verification lifecycle:
\begin{enumerate}
  \item \textbf{Intent parsing}: translate natural-language intent to formal policy
  \item \textbf{Feasibility checking}: can this intent be implemented in the current network topology?
  \item \textbf{Pre-deployment formal verification}: does the planned NF configuration satisfy the intent?
  \item \textbf{Post-deployment monitoring}: does the running NF chain continue to satisfy the intent over time?
\end{enumerate}

\subsubsection{Classification as NF Validation}

This is \emph{holistic NF lifecycle validation} --- arguably the most complete vision of NF correctness in our survey. Yaksha-Prashna currently occupies only Phase 3 (pre-deployment verification). This paper shows what the complete picture looks like and identifies where YP fits in the ecosystem.

\subsubsection{Technical Mechanism}

\paragraph{Phase 1 --- Intent Parsing.}
Operator intents are translated from natural language (or a constrained intent language like NEMO or ITDM) to formal network policies. The translation uses NLP techniques (dependency parsing, entity recognition for network objects like VLANs, services, NF types) combined with a formal policy language.

\paragraph{Phase 2 --- Feasibility Checking.}
Given the translated policy and the current network topology, check whether the policy \emph{can} be implemented: does the topology have enough NF instances, are the required NF types available, is bandwidth sufficient? This is a constraint satisfaction problem over the topology graph.

\paragraph{Phase 3 --- Pre-deployment Formal Verification.}
Given the planned NF configuration (which NFs, in which order, with which parameters), formally verify that the configuration satisfies the policy. This uses model checking or SMT-based verification. \textbf{This is where Yaksha-Prashna fits.}

\paragraph{Phase 4 --- Post-deployment Monitoring.}
After deployment, continuously monitor the running NF chain against the policy. This uses runtime telemetry (packet traces, flow statistics, NF state snapshots) and checks for policy violations. Violations trigger automated remediation (NF restart, configuration update, scaling).

\textbf{Tools}: NLP (for intent parsing), SMT/model checking (for Phase 3 verification), streaming analytics (for Phase 4 monitoring).

\begin{analogybox}
\textbf{YP as Phase 3 of the Full-Lifecycle framework}

The Full-Lifecycle paper provides the \emph{ecosystem context} for Yaksha-Prashna:

\begin{center}
\begin{tabular}{lll}
\toprule
\textbf{Lifecycle Phase} & \textbf{Full-Lifecycle Paper} & \textbf{Yaksha-Prashna} \\
\midrule
Intent parsing & NLP $\to$ formal policy & (out of scope) \\
Feasibility & Topology constraint check & (out of scope) \\
Pre-deployment & SMT/model checking & \textbf{This is YP's domain} \\
Post-deployment & Runtime monitoring & (out of scope) \\
\bottomrule
\end{tabular}
\end{center}

YP is the \textbf{formal verification engine} for Phase 3 of this lifecycle. The Full-Lifecycle paper provides:
\begin{enumerate}
  \item A framework positioning YP within a larger ecosystem (not ``just a verification tool'')
  \item A blueprint for what Phases 1, 2, 4 should look like if YP is to be part of a complete solution
  \item The idea of \emph{intent-to-policy translation} as input to YP's Prolog assertions
\end{enumerate}

What is particularly interesting: the paper shows that \emph{policy intent} (Phase 1 output) and \emph{behavioral assertions} (YP's input) are the same thing --- just expressed differently. If Phase 1's NLP can translate operator intents into YP's Prolog assertion format, YP becomes immediately accessible to non-expert operators.
\end{analogybox}

\begin{strategybox}
\textbf{Strategy T --- Intent-to-Prolog LLM bridge}:\\
Implement a natural-language-to-Prolog assertion translator for YP:
\begin{enumerate}
  \item Operator writes: ``All traffic from external interfaces must be inspected by the IDS before reaching the database''
  \item LLM (GPT-4 or Gemini) translates to Prolog:
  \begin{lstlisting}[language=Prolog]
chain_property(Chain) :-
    processes_before(Chain, ids, database_nf),
    forall(packet(P, external_interface), 
           reaches(P, ids)).
  \end{lstlisting}
  \item YP formally verifies this assertion against the compiled eBPF bytecode
  \item If verification fails: YP generates a counterexample packet + chain path
\end{enumerate}
This makes YP \textbf{usable by operators} who cannot write Prolog, dramatically expanding the user base.

\textbf{Strategy U --- Integrate YP into a full-lifecycle pipeline}:\\
Position YP as the core of a Full-Lifecycle NF verification system:
\begin{itemize}
  \item Phase 1: LLM intent parser $\to$ Prolog assertions (new component)
  \item Phase 2: Topology feasibility checker (new component, based on YP's chain graph)
  \item Phase 3: YP (existing, enhanced)
  \item Phase 4: eBPF-based runtime monitoring that re-uses YP's Prolog queries as streaming filters (new component)
\end{itemize}

\textbf{Strategy V --- Post-deployment monitoring using YP's analysis as streaming filters}:\\
YP's Prolog assertions are static properties. Compile them into streaming packet filters (using BPF program generation from Prolog rules) and deploy alongside the NF chain for continuous runtime verification. This closes the Phase 3 $\to$ Phase 4 loop.
\end{strategybox}

\begin{limitbox}
\begin{itemize}
  \item The paper is \textbf{architectural/conceptual}: it describes the lifecycle framework without implementing all phases.
  \item \textbf{Intent parsing accuracy}: NLP translation of complex intents to formal policies is error-prone; misinterpretation is a safety risk.
  \item Phase 4 monitoring incurs \textbf{runtime overhead}: continuous packet-level monitoring is expensive for high-speed NFs.
  \item The paper is \textbf{agnostic to NF implementation}: it does not specify how Phase 3 formal verification works --- just that it should. YP fills this gap.
\end{itemize}
\end{limitbox}

% ------------------------------------------------------------
\subsection{Synthesis: NetSMC + Full-Lifecycle = YP's Future}
% ------------------------------------------------------------

\begin{insightbox}
NetSMC and Full-Lifecycle Intent together define the two most important growth directions for Yaksha-Prashna:

\begin{enumerate}
  \item \textbf{Temporal reasoning} (from NetSMC): YP needs LTL-style temporal operators to express multi-packet, multi-hop chain properties. NetSMC's algorithms (containment optimization, synchronized product) directly inform the implementation.
  \item \textbf{Ecosystem integration} (from Full-Lifecycle): YP should not be a standalone tool but a formal verification \emph{module} in a complete intent-driven NF lifecycle management system. The Full-Lifecycle paper provides the blueprint.
\end{enumerate}

Combined, these two directions transform YP from a ``batch verification tool'' into a \textbf{live, intent-driven, temporally-aware NF correctness enforcement system}.
\end{insightbox}

\paragraph{The Research Trajectory.}
\begin{enumerate}
  \item \textbf{Short-term}: Add LTL temporal operators to YP's Prolog engine (from NetSMC).
  \item \textbf{Medium-term}: Implement intent-to-Prolog LLM bridge (from Full-Lifecycle Phase 1).
  \item \textbf{Long-term}: Compile YP's Prolog assertions into eBPF streaming monitors for Phase 4 (from Full-Lifecycle's post-deployment vision).
\end{enumerate}


% ============================================================
\section{Cross-Paper Strategy Matrix}
\label{sec:matrix}
% ============================================================

The following table synthesizes the research strategies extracted from all analyzed papers, mapping each to the specific YP limitation it addresses.

\begin{longtable}{p{3.5cm} p{4cm} p{4cm} p{3cm}}
\toprule
\textbf{Paper} & \textbf{Key Technique} & \textbf{YP Improvement} & \textbf{Addresses Limitation} \\
\midrule
\endhead
Heimdall & angr symbolic execution + Z3 equivalence & Symbolic enrichment layer; ITE map encoding & Path explosion; Map aliasing \\
\midrule
VigNAT & KLEE + VeriFast: DS split into formal + stateless & Split YP's map analysis: formal model DS ops + SE for logic & Stateful field limit; Map composition \\
\midrule
Vigor & Push-button verification framework & Design YP as a framework; Library abstraction & Usability; Integration \\
\midrule
BeePL & Type-preserving eBPF DSL compiler & Infer type information from verified bytecode; Type-aware CFG-NC & Limited field tracking \\
\midrule
AppNet & Z3 chain equivalence for ANF & Formal chain composition operator for YP; Z3 for chain equivalence & NF chain verification \\
\midrule
eNetSTL & Library contracts + verifier metadata & Treat kfunc annotations as trusted behavioral contracts in YP & Library calls; Verification coverage \\
\midrule
Demystifying eBPF & $\omega/\rho$ performance model & Add performance assertions to Prolog query vocabulary & Performance validation (new capability) \\
\midrule
FlowMage & GPT-4 field access extraction & YP does this formally; Add flow affinity queries to Prolog & Query vocabulary; Positioning \\
\midrule
NetSMC & Custom LTL model checker for NF chains & Add temporal operators to YP's Prolog; Adopt containment optimization & Temporal queries; Chain verification \\
\midrule
Full-Lifecycle Intent & IBN lifecycle (pre- + post-deploy) & Position YP as pre-deploy phase; Add intent-to-Prolog LLM translation & Ecosystem positioning; Lifecycle coverage \\
\bottomrule
\caption{Cross-paper strategy matrix: papers $\to$ YP improvements}
\label{tab:strategy-matrix}
\end{longtable}

% ============================================================
\section{Synthesis: Proposed Research Directions}
\label{sec:synthesis}
% ============================================================

Based on the cross-paper analysis, we propose the following concrete research directions for Yaksha-Prashna:

\subsection{Direction 1: Temporal Query Extensions (from NetSMC)}

Extend YP's Prolog query engine with \textbf{LTL-style temporal operators}:
\begin{itemize}
  \item $\mathbf{G}(P)$: ``Property $P$ holds for all packets in all executions''
  \item $\mathbf{F}(P)$: ``Property $P$ holds for some packet in some execution''
  \item $P\ \mathbf{U}\ Q$: ``$P$ holds until $Q$ is satisfied'' (e.g., ``firewall always sees packet before load balancer'')
\end{itemize}
Inspired by NetSMC's LTL model checker, but implemented as Prolog meta-predicates over YP's CFG-NC trace semantics.

\subsection{Direction 2: BPF Map Composition Algebra (from VigNAT + AppNet)}

Formalize BPF map composition across NF chains using:
\begin{itemize}
  \item \textbf{Ghost maps} (inspired by VigNAT/Klint): abstract BPF maps as uninterpreted SMT functions
  \item \textbf{Symbolic state transformers} (inspired by AppNet): model each NF as a function $(packet, state) \to (result, state')$
  \item Chain composition: sequential composition of state transformers, verified for equivalence or property satisfaction via Z3
\end{itemize}

\subsection{Direction 3: Type-Aware Analysis (from BeePL)}

Recover type information from eBPF verified bytecode to enrich CFG-NC:
\begin{itemize}
  \item The eBPF verifier propagates register types (PTR\_TO\_MAP\_VALUE, PTR\_TO\_PACKET, etc.)
  \item YP can read the verifier's type annotations from BTF (BPF Type Format) metadata embedded in the ELF
  \item Use these types to automatically expand the set of tracked fields beyond the current 6
\end{itemize}

\subsection{Direction 4: Library Contract Integration (from eNetSTL)}

Build a \textbf{YP annotation library}: for common eBPF library functions (Katran, Cilium helpers, eNetSTL kfuncs), pre-compute behavioral contracts and cache them. When YP encounters a call to a known library function, substitute the pre-computed contract instead of re-analyzing the implementation. This is analogous to VigNAT's ``trusted data structure library'' approach.

\subsection{Direction 5: Performance Assertion Layer (from Demystifying eBPF)}

Integrate the $\omega/\rho$ performance model into YP's static analysis:
\begin{itemize}
  \item At static time, classify each NF by hook type (XDP native $<$ TC $<$ SK\_SKB)
  \item Count per-packet operations (map lookups, helper calls, conditional branches)
  \item Estimate $\rho$ (overhead ratio) and check against $\omega$ (offload ratio) for the expected workload
  \item Report performance assertion violations: ``This NF's design will hurt performance for workloads with $\omega < \rho - 1$''
\end{itemize}

\subsection{Direction 6: Intent-to-Prolog LLM Bridge (from FlowMage + Full-Lifecycle)}

Leverage LLMs (as FlowMage does, but deterministically verified by YP) to:
\begin{itemize}
  \item Translate operator \emph{intents} (natural language) into YP Prolog assertion clauses
  \item Use YP to formally verify the generated assertions against the NF bytecode
  \item Close the loop: if verification fails, LLM generates a repair suggestion
\end{itemize}
This positions YP as the formal backbone of an intent-driven NF lifecycle (as the Full-Lifecycle paper envisions).

% ============================================================
\section{Conclusion}
\label{sec:conclusion}
% ============================================================

This research strategy document has analyzed 9 selected papers from a curated survey of 89 works in NF validation and network verification. Through detailed analysis of each paper's technical mechanism, analogy to Yaksha-Prashna, and specific extractable research strategies, we have identified 6 concrete research directions for improving Yaksha-Prashna.

The central finding of this survey is that the research community has developed powerful individual techniques --- symbolic execution (Heimdall), separation logic (VigNAT), push-button verification (Vigor), correct-by-construction type systems (BeePL), chain equivalence checking (AppNet, NetSMC), library contracts (eNetSTL), performance modeling (Demystifying eBPF), LLM-assisted analysis (FlowMage), and intent-driven lifecycle management (Full-Lifecycle) --- but none of these approaches addresses the full stack of requirements that Yaksha-Prashna is uniquely positioned to handle: \emph{bytecode-level, source-free, multi-program chain behavioral verification at deployment scale}.

Our research agenda is therefore not to replace YP with one of these approaches, but to selectively incorporate the most valuable technical insights from each into YP's existing framework, strengthening it incrementally toward a complete NF validation solution.

% ============================================================
\end{document}
